\lyseis{Β1.4}{Άθροισμα και διαφορά τμημάτων}
\sol{1} $\mathrm{ΑΓ}=3+4=7$, $\mathrm{ΑΔ}=7+2{,}5=9{,}5$, $\mathrm{ΒΔ}=9{,}5-3=6{,}5$.
\sol{2} $\mathrm{ΒΓ}=9-4{,}5=4{,}5$.
\sol{3} $5+3{,}5+6=14{,}5$.
\sol{4} α) $15{,}4$ \quad β) $18{,}5$ \quad γ) $11{,}2$ \quad δ) $24$.
\sol{5} $\mathrm{ΑΒ}=2\mathrm{ΒΓ}$ και $\mathrm{ΑΓ}=\mathrm{ΑΒ}+\mathrm{ΒΓ}=3\mathrm{ΒΓ}=12$, άρα $\mathrm{ΒΓ}=4$ και $\mathrm{ΑΒ}=8$.
\sol{6} α) Σωστό. β) Σωστό. γ) Λάθος: $4\alpha$. δ) Λάθος: $\mathrm{ΒΓ}=\mathrm{ΑΓ}-\mathrm{ΑΒ}$.
\sol{7} Περίμετρος $15{,}4$ μ. Αν μένει ανοιχτό ένα μέτρο, χρειάζονται $14{,}4$ μ. δίχτυ, δηλαδή $14{,}4\cdot3=43{,}2$ €.
\sol{8} Λιγότερο από $14$ μ., γιατί το ευθύγραμμο τμήμα είναι ο συντομότερος δρόμος.
\sol{9} $350+480=830$ μ.
\sol{10} $250+180+320=750$ μ. Περπάτησε $750-600=150$ μ. περισσότερο.
\sol{11} $12+9+14+10+8=53$ μ.
\sol{12} Αν είχαμε νήμα τεντωμένο ανάμεσα στα σημεία, κάθε στροφή θα το μεγάλωνε. Έτσι το ευθύγραμμο είναι το μικρότερο.
\sol{13} Το $\mathrm{ΑΓ}$ είναι μικρότερο από $5+7=12$ και μεγαλύτερο από $7-5=2$: $2<\mathrm{ΑΓ}<12$.
\sol{14} Η πτυχή είναι στο μέσο του $\mathrm{ΑΒ}$, δηλαδή στα $5$ μ. από κάθε άκρο.
\sol{15} Ένα ορθογώνιο $2\times6$ και ένα $3\times5$ έχουν και τα δύο περίμετρο $16$ αλλά δεν είναι ίσα.
\sol{16} Ελεύθερη απάντηση.
\sol{17} Ελεύθερη απάντηση. Το μήκος της τεθλασμένης είναι μεγαλύτερο από την απόσταση των άκρων της.
